% ---------------------------------------------------------------- chapter quotes
% One quotation opens each chapter: \chapterquote{key}{quotation}{author}{source}. The key is the
% chapter's label (I.1, II.2, A) or, for an unlabelled chapter, its title. A chapter without an entry
% opens without one.
%
% Every quotation here was checked against its source text (most of them against a scan or the
% full text of the original edition) before it went in. Translations are marked: where an
% established English translation was checked, it is named; otherwise the original was checked
% and the English is ours. Do not add a quotation that has not been checked the same way: many
% famous lines are misquoted or misattributed. (Example: Wiener's often-quoted "a method of
% controlling a system by reinserting into it the results of its past performance" is not the
% wording of the 1950 first edition, which is quoted below.)

% ----- front matter
\chapterquote{0}
  {A Governor is a part of a machine by means of which the velocity of the machine is kept nearly
   uniform, notwithstanding variations in the driving-power or the resistance.}
  {James Clerk Maxwell}
  {\emph{On Governors}, Proceedings of the Royal Society 16 (1868), first sentence}

\chapterquote{Prelude: a drone in the wind}
  {The alteration of motion is ever proportional to the motive force impressed; and is made in
   the direction of the right line in which that force is impressed.}
  {Isaac Newton}
  {\emph{Principia} (1687), Law II, in the translation of Andrew Motte (1729)}

% ----- Part I
\chapterquote{I.1}
  {\ldots feedback is the control of a system by reinserting into the system the results of its
   performance.}
  {Norbert Wiener}
  {\emph{The Human Use of Human Beings} (1950), p.~71}

\chapterquote{I.2}
  {Ut tensio sic vis; That is, The Power of any Spring is in the same proportion with the Tension
   thereof.}
  {Robert Hooke}
  {\emph{De Potentia Restitutiva, or of Spring} (1678)}

\chapterquote{I.3}
  {To every action there is always opposed an equal reaction: or the mutual actions of two bodies
   upon each other are always equal, and directed to contrary parts.}
  {Isaac Newton}
  {\emph{Principia} (1687), Law III, in the translation of Andrew Motte (1729)}

\chapterquote{I.4}
  {The harmonic oscillator, which we are about to study, has close analogs in many other fields;
   although we start with a mechanical example of a weight on a spring, \ldots\ we are really
   studying a certain differential equation.}
  {Richard Feynman}
  {\emph{The Feynman Lectures on Physics}, vol.~I (1963), §21-1}

\chapterquote{I.5}
  {And as every present state of a simple substance is naturally a consequence of its preceding
   state, in such a way that its present is big with its future.}
  {Gottfried Wilhelm Leibniz}
  {\emph{Monadology} (1714), §22, in the translation of Robert Latta (1898)}

\chapterquote{I.6}
  {For a successful technology, reality must take precedence over public relations, for nature
   cannot be fooled.}
  {Richard Feynman}
  {Report of the Presidential Commission on the Space Shuttle \emph{Challenger} Accident (1986),
   Appendix F}

\chapterquote{I.7}
  {Nothing takes place suddenly, and it is one of my great and best confirmed maxims that nature
   never makes leaps.}
  {Gottfried Wilhelm Leibniz}
  {\emph{New Essays on Human Understanding} (written 1704), Preface, in the translation of
   P.~Remnant and J.~Bennett (1981)}

\chapterquote{I.8}
  {I often say that when you can measure what you are speaking about, and express it in numbers,
   you know something about it; but when you cannot measure it, when you cannot express it in
   numbers, your knowledge is of a meagre and unsatisfactory kind.}
  {William Thomson, Lord Kelvin}
  {\emph{Electrical Units of Measurement}, lecture to the Institution of Civil Engineers, 3~May 1883}

\chapterquote{I.9}
  {If a function $f(t)$ contains no frequencies higher than $W$ cps, it is completely determined
   by giving its ordinates at a series of points spaced $1/2W$ seconds apart.}
  {Claude Shannon}
  {\emph{Communication in the Presence of Noise}, Proceedings of the IRE 37 (1949), Theorem~1}

\chapterquote{I.10}
  {On the other hand, under certain conditions of delay, etc., a feedback that is too brusque will
   make the rudder overshoot, and will be followed by a feedback in the other direction, which
   makes the rudder overshoot still more, until the steering mechanism goes into a wild
   oscillation or hunting, and breaks down completely.}
  {Norbert Wiener}
  {\emph{Cybernetics} (1948), Introduction}

\chapterquote{I.11}
  {\ldots big whirls have little whirls that feed on their velocity, and little whirls have lesser
   whirls and so on to viscosity\,---\,in the molecular sense.}
  {Lewis Fry Richardson}
  {\emph{Weather Prediction by Numerical Process} (1922), p.~66}

\chapterquote{I.12}
  {Give me a place to stand, and I will move the earth.}
  {Archimedes}
  {as reported by Pappus of Alexandria, \emph{Collection}, Book VIII (c.~AD~340); translated from
   the Greek}

\chapterquote{I.13}
  {Science is built up of facts, as a house is built of stones; but an accumulation of facts is no
   more a science than a heap of stones is a house.}
  {Henri Poincaré}
  {\emph{La Science et l'hypothèse} (1902), chapter~IX; translated from the French}

\chapterquote{I.14}
  {However, an excessive feedback is likely to be as serious a handicap to organized activity as a
   defective feedback.}
  {Norbert Wiener}
  {\emph{Cybernetics} (1948), Introduction}

% ----- Part II
\chapterquote{II.1}
  {We ought then to regard the present state of the universe as the effect of its anterior state
   and as the cause of the one which is to follow.}
  {Pierre-Simon Laplace}
  {\emph{A Philosophical Essay on Probabilities} (1814), in the translation of F.~W.~Truscott and
   F.~L.~Emory (1902)}

\chapterquote{II.2}
  {Nothing at all takes place in the world in which some rule of maximum or minimum does not
   appear.}
  {Leonhard Euler}
  {\emph{Methodus inveniendi lineas curvas} (1744), Appendix~I; translated from the Latin}

\chapterquote{II.3}
  {As far as the laws of mathematics refer to reality, they are not certain; and as far as they
   are certain, they do not refer to reality.}
  {Albert Einstein}
  {\emph{Geometry and Experience}, address to the Prussian Academy of Sciences (1921), in the
   translation of G.~B.~Jeffery and W.~Perrett (1922)}

\chapterquote{II.4}
  {This principle can be extended without difficulty to observations of unequal accuracy.
   \ldots\ the most probable system of values of the quantities \ldots\ will be that in which
   \ldots\ the sum of the squares of the differences between the actually observed and computed
   values multiplied by numbers that measure the degree of precision, is a minimum.}
  {Carl Friedrich Gauss}
  {\emph{Theoria motus corporum coelestium} (1809), art.~179, in the translation of C.~H.~Davis
   (1857)}

\chapterquote{II.5}
  {Probability is relative, in part to this ignorance, in part to our knowledge.}
  {Pierre-Simon Laplace}
  {\emph{A Philosophical Essay on Probabilities} (1814), in the translation of F.~W.~Truscott and
   F.~L.~Emory (1902)}

\chapterquote{II.6}
  {The first principle is that you must not fool yourself---and you are the easiest person to fool.}
  {Richard Feynman}
  {\emph{Cargo Cult Science}, Caltech commencement address (1974)}

\chapterquote{II.7}
  {Hitherto I have not been able to discover the cause of those properties of gravity from
   ph\ae nomena, and I frame no hypotheses.}
  {Isaac Newton}
  {\emph{Principia}, General Scholium (1713), in the translation of Andrew Motte (1729)}

\chapterquote{II.8}
  {Every body perseveres in its state of rest, or of uniform motion in a right line, unless it is
   compelled to change that state by forces impressed thereon.}
  {Isaac Newton}
  {\emph{Principia} (1687), Law I, in the translation of Andrew Motte (1729)}

\chapterquote{II.9}
  {If, for instance, I say, \enquote{That train arrives here at 7 o'clock}, I mean something like
   this: \enquote{The pointing of the small hand of my watch to 7 and the arrival of the train are
   simultaneous events.}}
  {Albert Einstein}
  {\emph{On the Electrodynamics of Moving Bodies} (1905), §1, in the translation of W.~Perrett and
   G.~B.~Jeffery (1923)}

\chapterquote{II.10}
  {One geometry cannot be more true than another; it can only be more convenient.}
  {Henri Poincaré}
  {\emph{La Science et l'hypothèse} (1902); translated from the French}

\chapterquote{II.11}
  {No figures will be found in this work. The methods I set out in it require neither
   constructions nor geometrical or mechanical reasoning, but only algebraic operations, subject to
   a regular and uniform procedure.}
  {Joseph-Louis Lagrange}
  {\emph{Méchanique analitique} (1788), Avertissement; translated from the French}

\chapterquote{II.12}
  {Whenever some change occurs in Nature, the quantity of action necessary for that change is the
   least possible.}
  {Pierre-Louis Moreau de Maupertuis}
  {\emph{Les Loix du mouvement et du repos déduites d'un principe metaphysique} (1746); translated
   from the French}

\chapterquote{II.13}
  {\ldots\ it is exceedingly important to shoot the missile, not at the target, but in such a way
   that missile and target may come together in space at some time in the future.}
  {Norbert Wiener}
  {\emph{Cybernetics} (1948), Introduction}

\chapterquote{II.14}
  {The present is never our end. The past and the present are our means; the future alone is our
   end.}
  {Blaise Pascal}
  {\emph{Pensées} (1670), §172, in the translation of W.~F.~Trotter}

\chapterquote{II.15}
  {The second, to divide each of the difficulties I would examine into as many parts as possible,
   and as would be required to solve them best.}
  {René Descartes}
  {\emph{Discours de la méthode} (1637), part~II, the second rule; translated from the French}

\chapterquote{II.16}
  {After spending a lot of time trying to estimate them by pure combinatorial calculations, I
   wondered whether a more practical method than \enquote{abstract thinking} might not be to lay it
   out say one hundred times and simply observe and count the number of successful plays.}
  {Stanisław Ulam}
  {on the chances of a game of solitaire, 1946; remarks of 1983 quoted by R.~Eckhardt, \emph{Los
   Alamos Science} 15 (1987)}

\chapterquote{II.17}
  {The motion of a system of material points, connected with one another in any way whatever,
   whose motions are at the same time bound by any external restrictions whatever, takes place at
   every instant in the greatest possible agreement with the free motion, or under the least
   possible constraint.}
  {Carl Friedrich Gauss}
  {\emph{Über ein neues allgemeines Grundgesetz der Mechanik}, Crelle's Journal 4 (1829); translated
   from the German}

\chapterquote{II.18}
  {If, by altering the adjustments of the machine, its governing power is continually increased,
   there is generally a limit at which the disturbance, instead of subsiding more rapidly, becomes
   an oscillating and jerking motion, increasing in violence till it reaches the limit of action
   of the governor.}
  {James Clerk Maxwell}
  {\emph{On Governors}, Proceedings of the Royal Society 16 (1868)}

\chapterquote{II.19}
  {Whenever therefore a bird pursues his course for some time without working his wings we must
   conclude either (1) that the course is not horizontal, (2) that the wind is not horizontal, or
   (3) that the wind is not uniform.}
  {John William Strutt, Lord Rayleigh}
  {\emph{The Soaring of Birds}, Nature 27 (1883)}

\chapterquote{II.20}
  {The Analytical Engine has no pretensions whatever to \emph{originate} any thing. It can do
   whatever \emph{we know how to order it} to perform.}
  {Ada Lovelace}
  {Notes by the translator of L.~F. Menabrea's \emph{Sketch of the Analytical Engine}, Scientific
   Memoirs 3 (1843)}

\chapterquote{II.21}
  {We are to admit no more causes of natural things than such as are both true and sufficient to
   explain their appearances.}
  {Isaac Newton}
  {\emph{Principia}, Book~III, Rule~I, in the translation of Andrew Motte (1729)}

% ----- the mathematical toolbox
\chapterquote{The mathematical toolbox}
  {Philosophy is written in this grand book that stands always open before our eyes (I mean the
   universe), but it cannot be understood unless one first learns to understand the language and
   to know the characters in which it is written. It is written in the language of mathematics.}
  {Galileo Galilei}
  {\emph{Il Saggiatore} (1623); translated from the Italian}

\chapterquote{A}
  {Data aequatione quotcunque fluentes quantitates involvente, fluxiones invenire; et vice versa.
   \textup{(}Given an equation involving any number of fluent quantities, to find the fluxions,
   and vice versa.\textup{)}}
  {Isaac Newton}
  {the sentence hidden in an anagram in his \emph{Epistola posterior} to Leibniz (1676)}

\chapterquote{B}
  {This condition is mathematically equivalent to the condition that all the possible roots, and
   all the possible parts of the impossible roots, of a certain equation shall be negative.}
  {James Clerk Maxwell}
  {\emph{On Governors}, Proceedings of the Royal Society 16 (1868)}

\chapterquote{C}
  {But then the rigorous logic of the matter is not plain! Well, what of that? Shall I refuse my
   dinner because I do not fully understand the process of digestion? No, not if I am satisfied
   with the result.}
  {Oliver Heaviside}
  {\emph{Electromagnetic Theory}, vol.~II (1899), §225}

\chapterquote{D}
  {God made the integers; all else is the work of man.}
  {Leopold Kronecker}
  {at the Berlin meeting of natural scientists, 1886, as reported by Heinrich Weber (1893);
   translated from the German}

\chapterquote{E}
  {It is seen in this essay that the theory of probabilities is at bottom only common sense
   reduced to calculus.}
  {Pierre-Simon Laplace}
  {\emph{A Philosophical Essay on Probabilities} (1814), in the translation of F.~W.~Truscott and
   F.~L.~Emory (1902)}

\chapterquote{F}
  {The purpose of computing is insight, not numbers.}
  {Richard Hamming}
  {\emph{Numerical Methods for Scientists and Engineers} (1962), motto}

\chapterquote{G}
  {A very small cause, which escapes our notice, determines a considerable effect which we cannot
   fail to see, and then we say that this effect is due to chance.}
  {Henri Poincaré}
  {\emph{Science et méthode} (1908), Book~I, chapter~4; translated from the French}

\chapterquote{H}
  {I then and there felt the galvanic circuit of thought close; and the sparks which fell from it
   were the fundamental equations between i, j, k; exactly such as I have used them ever since.}
  {William Rowan Hamilton}
  {letter to his son Archibald, 5~August 1865, recalling 16~October 1843}

\chapterquote{I}
  {It will be seen that matrices (attending only to those of the same order) comport themselves
   as single quantities; they may be added, multiplied or compounded together, \&c.}
  {Arthur Cayley}
  {\emph{A Memoir on the Theory of Matrices}, Philosophical Transactions 148 (1858)}

\chapterquote{J}
  {One will add the various differential functions that must vanish by the conditions of the
   problem, after multiplying each of them by an undetermined coefficient; one will set the whole
   equal to zero, and one will thus have a differential equation to be treated as an ordinary
   equation \emph{de maximis et minimis}.}
  {Joseph-Louis Lagrange}
  {\emph{Méchanique analitique} (1788), p.~46; translated from the French}
